%\usepackage[a4paper,hmargin=3.5cm,vmargin=5.3cm,head=1.5cm,foot=3.2cm]{geometry}
\usepackage[
paperwidth=190mm,
paperheight=260mm,
top=0.98in,
bottom=0.98in,
left=0.98in,
right=0.98in,
headsep=0pt,
footskip=0.4in
]{geometry}

% If you use headers/footers later:
%\setlength{\headheight}{0pt}

\usepackage{url}          % simple URL typesetting

% https://tex.stackexchange.com/questions/47576/combining-ifxetex-and-ifluatex-with-the-logical-or-operation
\usepackage{ifxetex,ifluatex}
\newif\ifxetexorluatex
\ifxetex
  \xetexorluatextrue
\else
  \ifluatex
    \xetexorluatextrue
  \else
    \xetexorluatexfalse
  \fi
\fi

\ifxetexorluatex
  % To fix LaTeX command is already defined error, load amssymb before unicode-math
  % https://tex.stackexchange.com/a/547153/56364
  % \usepackage{stix}
  \usepackage{amssymb}
  \usepackage{amsthm}
  \usepackage{fontspec}
  \usepackage{unicode-math}

  \defaultfontfeatures{Scale=MatchLowercase,Ligatures=TeX}
  % \setmainfont{STIX Two Text}
  \setmainfont{Noto Serif}
  \setmathfont{texgyreschola-math.otf}
  \usepackage[scale=0.8]{noto-mono}
  \setmonofont{Noto Sans Mono}
  % small url fonts
  \renewcommand{\UrlFont}{\ttfamily\small}

  \ifxetex
    % https://gist.github.com/dlimpid/5454229
    % Use xeCJK in XeLaTeX
    \usepackage{xeCJK}
    \xeCJKsetup{%
      CJKspace=true,% % true이면 띄어쓰기 사용. 중국어, 일어는 필요 없을수도
      % CJKmath=true,%  % true면 math environment 안에서 CJK 글자 사용
      CJKecglue={}%   % Western과 CJK 사이의 공백 지정: {}로 간격을 없앰
    }
    \setCJKmainfont[
      Path=fonts/,
      Extension=.otf,
      BoldFont=NotoSerifKR-Bold,
      AutoFakeSlant,
      Ligatures=TeX]{NotoSerifKR-Regular.otf}
    \setCJKsansfont[
      Path=fonts/,
      Extension=.otf,
      BoldFont=NotoSerifKR-Bold,
      AutoFakeSlant,
      Ligatures=TeX]{NotoSerifKR-Regular.otf}
    \setCJKmonofont[
      Path=fonts/,
      Extension=.ttf,
      BoldFont=NotoSansMono-Bold,
      AutoFakeSlant,
      Ligatures=TeX]{NotoSansMono-Regular.ttf}
  \fi

  \ifluatex
    \usepackage[hangul]{luatexko}
    \setmainhangulfont[
      Path=fonts/,
      UprightFont=*-Regular,
      BoldFont=*-Bold,
      Extension=.otf,
      AutoFakeSlant,
      Ligatures=TeX]{NotoSerifKR}
    \hangulbyhangulfont=1
    \setmainfont{Noto Serif}
    \setmathfont{texgyreschola-math.otf}
    \usepackage[scale=0.8]{noto-mono}
    \setmonofont{Noto Sans Mono}
  \fi
\else
  % pdfLaTeX fallback for broad compatibility with MiKTeX/TeX Live installations
  % https://t.co/PgBfFh9rTU?amp=1
  \usepackage{amsmath,amsthm,amsfonts}
  \usepackage[T1]{fontenc}
  \usepackage[utf8]{inputenc}
  \usepackage[finemath]{kotex} % Use kotex in PDFLaTeX
  % \usepackage{CJKutf8}
  \usepackage{textcomp} % provide euro and other symbols

  % use nanum font in Korean
  \usepackage{dhucs-nanumfont}
  % otherwise, use Noto fonts
  \usepackage{noto}
\fi

\ifxetex
  \XeTeXinputnormalization=1
\fi

\renewcommand{\floatpagefraction}{0.1} % one figure in one page

\usepackage[english]{babel}    % babel set language

% NeurIPS-style text fonts and citations, while keeping the thesis layout.
% The official NeurIPS style uses Times/Helvetica and natbib.
\usepackage{times}

% Citation style switch:
%   \numericcitationsfalse -> author-year citations, e.g. (Farronato et al., 2025)
%   \numericcitationstrue  -> numeric citations, e.g. [1]
\newif\ifnumericcitations
\numericcitationstrue

\ifnumericcitations
  \usepackage[numbers,square]{natbib}
\else
  \usepackage[authoryear,round]{natbib}
\fi
\newcommand{\textcite}{\citet}
\newcommand{\parencite}{\citep}
\usepackage{csquotes}

\usepackage{setspace}   %use this package to set linespacing as desired
\doublespacing
\emergencystretch=3em
\usepackage{afterpage}    % separate page for floats and tables
\usepackage{titlefoot}    % keywords on footnote
\usepackage[explicit]{titlesec}  %title control and formatting
\usepackage[titles]{tocloft}  %table of contents control and formatting
\usepackage[page]{appendix}  %for appendices
\usepackage[normalem]{ulem}  %for underlined section titles

\usepackage{graphicx}  %for images and plots
\usepackage{pdfpages}  %for inserting the signed approval page
\usepackage[figuresright]{rotating}  %for legacy rotated floats
\usepackage{pdflscape}  % rotate landscape pages in PDF viewers
\usepackage{adjustbox}    % fit table to page width
\usepackage{algorithm}      % algorithm
\usepackage{algpseudocode}  % algorithm

\newenvironment{landscapefigure}[1][p]{%
  \begin{landscape}%
  \begin{figure}[#1]%
}{%
  \end{figure}%
  \end{landscape}%
}

\newenvironment{landscapetable}[1][p]{%
  \begin{landscape}%
  \begin{table}[#1]%
}{%
  \end{table}%
  \end{landscape}%
}

\usepackage{booktabs}     % proessional-quality tables
\usepackage{longtable}    % multipage tables
\usepackage{multicol} 		% use multicol in table
\usepackage{multirow} 		% use multirow in table

\usepackage{xcolor}			  % color fonts
\usepackage[super]{nth}		% th
\usepackage{siunitx}		  % SI unit


%\usepackage{fontspec}
\usepackage{kotex} % Korean support in XeLaTeX


\usepackage[
  bookmarks=true,
  %colorlinks=true,
  %citecolor=blue,
  %linkcolor=blue,
  urlcolor=blue
]{hyperref}     % hyperlink; load after natbib so citation links work
\usepackage{bookmark}
\usepackage[noabbrev,nameinlink]{cleveref}		% for multiple cross-reference
\usepackage[intoc]{nomencl}  % nomenclature
\makenomenclature
\setlength{\nomitemsep}{-\parsep}
\setlength{\nomlabelwidth}{1.7cm}

% Theorem-like environments
\theoremstyle{plain}
\newtheorem{theorem}{Theorem}[chapter]
\newtheorem{lemma}[theorem]{Lemma}
\newtheorem{proposition}[theorem]{Proposition}
\theoremstyle{definition}
\newtheorem{assumption}[theorem]{Assumption}
\renewcommand{\thetheorem}{\arabic{chapter}.\arabic{theorem}}
\crefname{theorem}{theorem}{theorems}
\Crefname{theorem}{Theorem}{Theorems}
\crefname{lemma}{lemma}{lemmas}
\Crefname{lemma}{Lemma}{Lemmas}
\crefname{proposition}{proposition}{propositions}
\Crefname{proposition}{Proposition}{Propositions}
\crefname{assumption}{assumption}{assumptions}
\Crefname{assumption}{Assumption}{Assumptions}

%%%%%%%%%%%%%%%%%%%%%%%%%%%%%%%%%%%
% Bibliography
%%%%%%%%%%%%%%%%%%%%%%%%%%%%%%%%%%%

\renewcommand{\bibname}{References}
\providecommand{\refname}{References}
\renewcommand{\refname}{References}
\AtBeginDocument{%
  \renewcommand{\bibname}{References}%
  \renewcommand{\refname}{References}%
}

% footnote for keywords in abstract
% https://tex.stackexchange.com/questions/30720/footnote-without-a-marker
\newcommand\blfootnote[1]{%
  \begingroup
  \renewcommand\thefootnote{}\footnote{#1}%
  \addtocounter{footnote}{-1}%
  \endgroup
}

% Symbol-marked footnotes for notes that should not look like numeric citations.
\newcommand\symbolfootnote[2][\ensuremath{\dagger}]{%
  \begingroup
  \renewcommand{\thefootnote}{#1}%
  \footnotemark
  \footnotetext{#2}%
  \endgroup
  \addtocounter{footnote}{-1}%
}
%%%%%%%%%%%%%%%%%%%%%%
% Start of Document
%%%%%%%%%%%%%%%%%%%%%%



% Format for Table of Contents
\renewcommand{\cftchapdotsep}{\cftdotsep}  %add dot separators
\renewcommand{\cftchapfont}{\bfseries}  %set title font weight
\renewcommand{\cftchappagefont}{}  %set page number font weight
\renewcommand{\cftchappresnum}{Chapter }
\renewcommand{\cftchapaftersnum}{:}
\renewcommand{\cftchapnumwidth}{7em}
\setlength{\cftbeforechapskip}{\baselineskip}
\setlength{\cftbeforesecskip}{4pt}
\setlength{\cftbeforesubsecskip}{4pt}
\setlength{\cftbeforesubsubsecskip}{4pt}
\setcounter{secnumdepth}{3}
\setcounter{tocdepth}{3}

%format title font size and position (this also applys to list of figures and list of tables)
\titleformat{\chapter}[display]
{\normalfont\bfseries\filcenter}{\chaptertitlename\ \thechapter}{0pt}{\MakeUppercase{#1}}

% Adjust chapter title formatting
\newcommand{\chapternamefont}{\scshape\Large}% Chapter name font
\newcommand{\chaptertitlefont}{\LARGE\bfseries}% Chapter title font
\titleformat{\chapter}[display]
{\normalfont\bfseries\chapternamefont\raggedright}{\MakeUppercase\chaptertitlefont\chaptertitlename\ \thechapter}{0pt}{\MakeUppercase{#1}}  %spacing between titles
\titlespacing*{\chapter}
  {0pt}{\topskip}{30pt}	%controls vertical margins on title

% Adjust section title formatting
\titleformat{\section}{\normalfont\bfseries}{\thesection}{1em}{#1}

% Adjust subsection title formatting
% \titleformat{\subsection}{\normalfont}{\uline{\thesubsection}}{0em}{\uline{\hspace{1em}#1}}
\titleformat{\subsection}{\normalfont\bfseries}{\thesubsection}{1em}{#1}

% Adjust subsubsection title formatting
\titleformat{\subsubsection}[hang]{\normalfont\bfseries}{\thesubsubsection}{1em}{#1}

% Change chapter numbers to Roman numerals (I, II, III...)
\renewcommand{\thechapter}{\Roman{chapter}}
% Ensure sections, figures, tables, and equations still use arabic chapter numbers (e.g., 2.1, not II.1)
\renewcommand{\thesection}{\arabic{chapter}.\arabic{section}}
\renewcommand{\thesubsubsection}{\thesubsection~\Alph{subsubsection}}
\makeatletter
\renewcommand{\p@subsubsection}{}
\renewcommand{\theHsubsubsection}{\arabic{chapter}.\arabic{section}.\arabic{subsection}.\arabic{subsubsection}}
\makeatother
\renewcommand{\thefigure}{\arabic{chapter}.\arabic{figure}}
\renewcommand{\thetable}{\arabic{chapter}.\arabic{table}}
\renewcommand{\theequation}{\arabic{chapter}.\arabic{equation}}

% Format page numbers as -1- in the footer
\usepackage{fancyhdr}
\fancypagestyle{plain}{
  \fancyhf{}
  \renewcommand{\headrulewidth}{0pt}
  \fancyfoot[C]{- \thepage\ -}
}
